\subsection{Tensor products and rank characters}
The following uses the ordinary tensor product of semimodules over the
commutative semiring $R_0$. Write $\operatorname{Pic}$ for the group of
isomorphism classes of tensor-invertible semimodules.

\begin{proposition}\label{prop:tensor-picard}
For a commutative unital ring $R$, including the zero ring,
\[
 \mathcal P(P,(M_p))\otimes_{R_0}\mathcal P(Q,(N_q))
 \simeq\mathcal P(P\times Q,(M_p\otimes_R N_q)).
\]
Zero labels are allowed. Let $\mathcal C_R$ be the commutative monoid of
isomorphism classes of nonempty connected finite $V(R)$-labelled posets,
with this product and unit the singleton labelled $R$. Then, as a
semiring and a ring respectively,
\[
 V(R_0)\simeq\N[\mathcal C_R],\qquad
 K_0(R_0)\simeq\Z[\mathcal C_R],\qquad
 \operatorname{Pic}(R_0)\simeq\operatorname{Pic}(R).
\]
If $\mathfrak p\in\Spec R$, let $r(M)$ be the rank of $M$ at
$\mathfrak p$. For every $s\in\mathbb C$ the formula
\[
 D_s^+(\mathcal P(P,(M_p)))
   =\sum_{\substack{p\in P\\r(M_p)>0}}r(M_p)^{-s}
\]
extends to a ring homomorphism $K_0(R_0)\to\mathbb C$.
For actual projectives, $D_{-1}^+\ge D_0^+$.
\end{proposition}
\begin{proof}
Choose a split free presentation of each $M_p$, with projector $a_p$
on $R^{n_p}$ and $n_p\ge1$; use $n_p=1$ and $a_p=0$ for a zero label.
The construction in the proof of Theorem~\ref{thm:main}(4) presents its
semimodule by a matrix with full occupied support block from $q$ to $p$
exactly when $p\le q$, and ring-amplitude matrix block diagonal with
blocks $a_p$. Thus the other occupied entries have old-zero amplitude.
Tensor the two split presentations. The free basis is the Cartesian
product of the original bases, and the projector is their Kronecker
product. Its support order is the product order, while its diagonal
ring block is $a_p\otimes b_q$, with image $M_p\otimes_R N_q$.
This proves the tensor formula, including zero labels and empty posets.
A product of two nonempty connected posets is connected: move along a
comparability path in one coordinate with the other fixed, and then in
the other coordinate. Direct sum is disjoint union by
Theorem~\ref{thm:Kzero}. These observations give the monoid semiring
and its ring completion.

Every tensor-invertible semimodule $X$ is dualizable, with its tensor
inverse $Y$. Write the coevaluation of $1$ as a finite sum
$\sum_i x_i\otimes y_i$. The triangle identity gives
$x=\sum_i x_i\operatorname{ev}(y_i\otimes x)$. Hence the coefficient
maps $x\mapsto\operatorname{ev}(y_i\otimes x)$ split the surjection
$R_0^n\to X$ defined by the $x_i$: $X$ is finite projective.
The same holds for $Y$. The tensor formula now forces their support
posets to be singletons, since their product is the singleton poset
of $R_0$. Their labels $M,N$ satisfy $M\otimes_R N\simeq R$.
Conversely such invertible ring modules give inverse singleton-labelled
semimodules. This proves the Picard comparison. These projectives have
a zero least group; they are not the ordinary ring modules obtained
by restriction along the ring reflection.

Local rank is a nonnegative integer, is additive on ring direct sums,
and satisfies $r(M\otimes_R N)=r(M)r(N)$. A product rank is positive
exactly when both ranks are positive. The tensor formula therefore
factors the displayed sum, while disjoint union makes it additive.
The unit has rank one, so group completion gives the asserted unital
ring homomorphism. Finally each retained positive rank is at least one,
proving the inequality on actual projectives.
\end{proof}
The restriction to positive ranks is necessary. A singleton with zero
label has total rank zero but one support vertex; counting that vertex
at $s=0$ would give a false inequality. At positive real $s$ its
unrestricted expression $0^{-s}$ is undefined. Even a nonzero module
may have rank zero at a chosen prime of a disconnected ring. No rank
character is asserted for the zero ring, whose prime spectrum is empty.
